\section{Discussion}

We addressed two key challenges that limit the study of test code evolution: reliable test-method history tracking and method-level production-to-test mapping. Our first finding (\textbf{RQ1}) is encouraging: history-generation tools developed for production methods generalize well to test methods without any threshold recalibration. In particular, \emph{HistoryFinder}~\cite{islam_historyfinder_2025} achieved the best overall performance in both accuracy and runtime. In contrast, our evaluation of production-to-test mapping (\textbf{RQ2}) showed that no approach consistently achieves both high precision and high recall. Thus, the choice should be guided by the downstream objective. For example, precision-oriented approaches suit studies requiring reliable links, whereas recall-oriented approaches better support applications prioritizing broader coverage.

These validated techniques enabled a more accurate investigation of test method evolution. We found (\textbf{RQ3}) that nearly one-third of mapped test methods are revised more often than the production methods they exercise, while the correlation between production and test revisions is generally weak. This suggests that much of test maintenance cannot be explained solely by production-test co-evolution and that test code should be treated as a first-class software artifact with its own maintenance dynamics rather than merely a passive consequence of production-code changes.

Our analysis of test smells provides one explanation for this additional maintenance effort. Test methods that evolve more than their corresponding production methods contain more test smells (\textbf{RQ4}). Surprisingly, the association is strongest in small test methods, contradicting Spadini et al.~\cite{Spadini:2018}. This difference likely stems from our more accurate method-history tracking, production-to-test method alignment, and a larger benchmark (110 vs. 10 projects). The stronger impact of test smells in smaller methods underscores the importance of evaluating software quality indicators at the method level rather than extrapolating conclusions from file/class-level analyses.

\textbf{For researchers}, our findings have several implications. The ability to track test method evolution alongside production changes may help more reliable studies of test obsolescence by precisely identifying when a test becomes outdated relative to the production change that invalidated it. It may also support the study of flaky tests caused by evolutionary mismatch. In addition, it can improve regression testing research by enabling fine-grained identification of which tests should be re-run or updated after specific production changes. For test repair, it can help pinpoint the exact divergence point and derive minimal corrective fixes. Finally, it enables more accurate long-term test maintenance modeling by supporting measurement of technical debt accumulation, smell persistence, and maintenance effort over time.

% \textbf{For researchers}, our findings have several implications. The ability to track test method evolution alongside production changes enables more reliable studies of test obsolescence by precisely identifying when a test becomes outdated relative to the production change that invalidated it; flaky tests caused by evolutionary mismatch by capturing how tests and production co-evolve over time; regression testing by enabling fine-grained identification of which tests should be re-run or updated after specific production changes; automated test repair by pinpointing the exact divergence point and supporting minimal corrective fixes; and long-term test maintenance modeling by enabling accurate measurement of technical debt accumulation, smell persistence, and maintenance effort over time. Our results also suggest that future production-test co-evolution studies should move beyond the traditional assumption~\cite{zaidman2011studying, sun_revisiting_2023} that simultaneous or temporally close changes necessarily indicate co-evolution, since many test methods evolve independently.

\textbf{For practitioners}, our results indicate that test maintenance represents a substantial portion of long-term maintenance effort and deserves no less engineering attention than production code. Organizations should therefore monitor the evolution of test code independently instead of assuming that production-code metrics adequately reflect testing effort. The strong relationship between test smells and highly evolving tests also suggests practical opportunities for maintenance prioritization~\cite{Khomh:2012}. 
